\section{Experimental Results}
\label{sec:results}

Every quantitative claim, table entry, and figure coordinate in this section traces deterministically to \TotalSealedRuns{} SHA-256-sealed v3 suites (\TotalSealedRecords{} simulation records + \BenchTotalBatches{} real executed Python batches, \TotalEscapedDefects{} cumulative mainline escapes).

\subsection{RQ1: Executable 64-Batch Multi-PR Causal Blame}
Table~\ref{tab:blame_microbench} and Figure~\ref{fig:blame_regime}(a) report attribution accuracy across all \BenchTotalBatches{} real executed failing Python batches. Simple heuristics perform poorly: \texttt{RandomRanker} achieves \BenchRandTopOnePct{}\% Top-1 accuracy, while \texttt{ChurnRiskRanker} scores \BenchChurnTopOnePct{}\% Top-1 because subtle single-line semantic defects are smaller than benign functional refactors and documentation updates.

% AUTO-GENERATED FROM v2_exp0_blame_microbench — DO NOT EDIT
\begin{table}[t]
\centering
\caption{Real Executed Python Multi-PR Blame Micro-Benchmark (48 Failing Batches: 24 Regressions + 24 True Pairwise Semantic Conflicts).}
\label{tab:blame_microbench}
\begin{tabular}{lccc}
\toprule
\textbf{Ranker} & \textbf{Top-1 Acc.} & \textbf{Top-2 Acc.} & \textbf{Latency} \\
\midrule
RandomRanker & 29.2\% & 58.3\% & $<0.1$~ms \\
ChurnRiskRanker & 2.1\% & 8.3\% & $<0.1$~ms \\
\textbf{ASTStackCausalRanker (Ours)} & \textbf{91.7\%} & \textbf{100.0\%} & \textbf{2.29~ms} \\
\bottomrule
\end{tabular}
\end{table}


The strong non-AST baseline \texttt{FileTracebackOverlapRanker} parses the identical Python stack traces and test module hints, lifting overall Top-1 accuracy to \BenchFileTopOnePct{}\% (\BenchFileTopTwoPct{}\% Top-2, MRR \BenchFileMRR{}) and achieving \BenchFileSemPct{}\% Top-1 on cross-file semantic conflicts where both conflicting modules appear on the stack. Crucially, however, \texttt{FileTracebackOverlapRanker} collapses to \BenchFileCofilePct{}\% Top-1 on the \BenchCofileBatches{} \emph{co-file distractor} batches: when a benign documentation or comment PR edits the same file with larger line churn, file-level overlap cannot separate syntactic noise from executable code changes. By filtering AST-invariant edits and expanding a 1-hop function call graph from failing tests, \texttt{ASTStackCausalRanker} achieves \BenchASTTopOnePct{}\% overall Top-1 accuracy, \BenchASTTopTwoPct{}\% Top-2 accuracy, \BenchASTCofilePct{}\% on co-file distractors, \BenchASTSemPct{}\% on semantic conflicts, and \BenchASTMRR{} MRR in just \BenchASTMs{}~ms per failure---confirming that lightweight static analysis provides high-precision, zero-LLM blame attribution.

% AUTO-GENERATED FROM v4_exp2_main_bench — DO NOT EDIT
\begin{table*}[t]
\centering
\footnotesize
\setlength{\tabcolsep}{4.2pt}
\caption{Main Benchmark Across 3 Real OSS Traces (450 Eval PRs, $5\times$ Velocity, $\phi=5\%$, $\rho=80\%$, 20 Paired Seeds $\times$ 3 Repos; 0 Cumulative Escapes. Throughput, P95, and Cost Show Mean [95\% Marginal Bootstrap CI]; Rej.\ Shows Mean Defective/Persistent/Flaky Counts).}
\label{tab:main_benchmark}
\begin{tabular}{lcccccc}
\toprule
\textbf{Strategy} & \textbf{Throughput (PR/h)} $\uparrow$ & \textbf{P95 Lat.\ (min)} $\downarrow$ & \textbf{CI Cost (m/PR)} $\downarrow$ & \textbf{Inv.\ (\%)} $\downarrow$ & \textbf{FR (\%)} $\downarrow$ & \textbf{Rej.\ (D/P/F)} \\
\midrule
SerialQueue & 3.75 [3.68, 3.81] & 1413.9 [1406.2, 1421.3] & 16.11 [15.80, 16.43] & 0.0\% & 4.36\% & 11.8/19.8/6.0 \\
StaticBatchBisect & 16.92 [16.32, 17.51] & 113.8 [105.6, 122.4] & 18.25 [17.35, 19.22] & 0.9\% & 2.01\% & 11.8/20.2/2.8 \\
SpeculativeLinear & 5.82 [5.57, 6.06] & 846.8 [818.6, 875.0] & 28.63 [27.44, 29.86] & 38.3\% & 2.12\% & 11.8/20.0/2.9 \\
FIFOOutcomeTree & 11.05 [10.59, 11.50] & 319.4 [304.7, 334.4] & 39.90 [38.35, 41.45] & 75.2\% & 1.78\% & 11.8/20.0/2.5 \\
DepAwareOutcomeTree & 11.75 [11.19, 12.29] & 283.5 [265.1, 302.0] & 36.95 [35.17, 38.87] & 71.5\% & 2.19\% & 11.8/20.0/3.0 \\
FlakeAwareBatch & 13.44 [12.79, 14.09] & 215.8 [197.5, 234.4] & 27.34 [25.66, 29.06] & 1.9\% & 0.12\% & 11.8/20.6/0.2 \\
\textbf{SSMQ} & \textbf{18.18 [17.66, 18.69]} & \textbf{81.0 [75.4, 86.7]} & \textbf{16.75 [16.02, 17.51]} & \textbf{2.0\%} & \textbf{1.17\%} & \textbf{11.8/20.2/1.6} \\
\bottomrule
\end{tabular}
\end{table*}


\subsection{RQ2: End-to-End Throughput, Tail Latency, Cost, and Scaling}
Table~\ref{tab:main_benchmark} summarizes aggregate performance across all three repositories at $5\times$ velocity, $\phi=5\%$ flakiness, and $\rho=80\%$ graph recall ($20\%$ latent off-graph conflicts):
\begin{itemize}
    \item \textbf{Vs.\ Static FIFO Batching}: SSMQ increases merge throughput by +\GainTputVsStaticPct{}\% (\SSMQThroughput{} vs.\ \StaticThroughput{}~PR/h), reduces P95 merge latency by \LatRedVsStaticPct{}\% (\SSMQPNinetyFiveLat{} vs.\ \StaticPNinetyFiveLat{}~min; Holm-adjusted Wilcoxon $p\HolmPLatVsStatic{}$, large effect Cliff's $|\delta|=\CliffsLatVsStatic{}$), lowers CI compute cost by \CostSaveVsStaticPct{}\% (\SSMQCICost{} vs.\ \StaticCICost{}~runner-min/PR, $p\HolmPCostVsStatic{}$), and cuts mean failure recovery rounds from \StaticBlameRounds{} to \SSMQBlameRounds{}.
    \item \textbf{Vs.\ EuroSys~'19 Outcome-Tree SubmitQueue}: Upgrading \texttt{UberSubmitQueue} to full pass/fail outcome-tree speculation nearly doubles throughput over \texttt{SpeculativeLinear} (\UberThroughput{} vs.\ \SpecLinThroughput{}~PR/h) and slashes P95 latency from \SpecLinPNinetyFiveLat{} to \UberPNinetyFiveLat{}~min. Yet because outcome-tree speculation still packs FIFO batches over overlapping modules, \UberInvalidRatePct{}\% of speculative runs are invalidated, burning \UberCICost{}~runner-min/PR. SSMQ delivers +\GainTputVsUberPct{}\% throughput, -\LatRedVsUberPct{}\% P95 latency ($p\HolmPLatVsUber{}$, $|\delta|=\CliffsLatVsUber{}$), and -\CostSaveVsUberPct{}\% CI cost savings.
    \item \textbf{Vs.\ Fixed Flaky Retry Batching}: \texttt{FlakeAwareBatch} nearly eliminates clean-PR false rejections (\FlakeBatchFalseRejPct{}\%) by retrying every failing batch once, but incurs severe tail inflation (\FlakeBatchPNinetyFiveLat{}~min P95, \FlakeBatchCICost{}~min/PR). SSMQ slashes P95 latency by \LatRedVsFlakePct{}\% and runner cost by \CostSaveVsFlakePct{}\% while keeping clean-PR false rejection at \SSMQFalseRejPct{}\% (vs.\ \StaticFalseRejPct{}\% for static batching, $p\HolmPFRVsStatic{}$).
\end{itemize}

% AUTO-GENERATED FROM v3_exp2_main_bench — DO NOT EDIT
\begin{figure*}[t]
\centering
\begin{tikzpicture}
\begin{groupplot}[
  group style={group size=2 by 1, horizontal sep=1.4cm},
  width=0.47\textwidth, height=4.8cm,
  symbolic x coords={Kubernetes,TypeScript,Next.js},
  xtick=data, ymin=0,
  ymajorgrids=true, grid style={dashed,gray!35},
  tick label style={font=\footnotesize},
  label style={font=\footnotesize},
  title style={font=\footnotesize\bfseries},
  enlarge x limits=0.22,
  error bars/y dir=both, error bars/y explicit,
  error bars/error bar style={line width=0.6pt, black!75},
]
\nextgroupplot[ybar=2pt, /pgf/bar width=7.5pt, title={(a) P95 Merge Latency (min) $\downarrow$}, ylabel={P95 Latency (min)}, ymax=470,
  legend to name=legPerRepo, legend columns=4, legend style={font=\footnotesize, draw=none, /tikz/every even column/.style={column sep=8pt}}]
\addplot[fill=ssmqGray!80, draw=ssmqGray!90!black] coordinates { (Kubernetes,92.70) +- (4.74,4.18) (TypeScript,167.20) +- (6.36,6.83) (Next.js,105.76) +- (6.93,6.97) };
\addlegendentry{StaticBatch}
\addplot[fill=ssmqOrange!80, draw=ssmqOrange!90!black] coordinates { (Kubernetes,323.59) +- (13.81,12.73) (TypeScript,411.67) +- (11.89,11.13) (Next.js,292.78) +- (15.49,16.92) };
\addlegendentry{UberTree}
\addplot[fill=ssmqPurple!80, draw=ssmqPurple!90!black] coordinates { (Kubernetes,181.50) +- (9.60,9.98) (TypeScript,331.92) +- (11.08,11.57) (Next.js,193.20) +- (14.71,14.46) };
\addlegendentry{FlakeBatch}
\addplot[fill=ssmqBlue!80, draw=ssmqBlue!90!black] coordinates { (Kubernetes,63.92) +- (4.82,4.64) (TypeScript,103.63) +- (8.10,8.54) (Next.js,70.86) +- (6.42,7.17) };
\addlegendentry{SSMQ (Ours)}
\nextgroupplot[ybar=2pt, /pgf/bar width=7.5pt, title={(b) CI Compute Cost (runner-min/PR) $\downarrow$}, ylabel={Runner-min / Merged PR}, ymax=60]
\addplot[fill=ssmqGray!80, draw=ssmqGray!90!black] coordinates { (Kubernetes,16.88) +- (0.51,0.47) (TypeScript,25.36) +- (0.83,0.86) (Next.js,17.09) +- (0.54,0.55) };
\addplot[fill=ssmqOrange!80, draw=ssmqOrange!90!black] coordinates { (Kubernetes,40.50) +- (1.67,1.74) (TypeScript,51.00) +- (1.60,1.56) (Next.js,38.77) +- (1.59,1.79) };
\addplot[fill=ssmqPurple!80, draw=ssmqPurple!90!black] coordinates { (Kubernetes,24.33) +- (0.86,0.88) (TypeScript,39.88) +- (1.48,1.40) (Next.js,25.56) +- (1.16,1.20) };
\addplot[fill=ssmqBlue!80, draw=ssmqBlue!90!black] coordinates { (Kubernetes,15.46) +- (0.65,0.68) (TypeScript,21.12) +- (0.92,0.99) (Next.js,14.58) +- (0.70,0.76) };
\end{groupplot}
\node at ($(group c1r1.north)!0.5!(group c2r1.north)+(0,0.85cm)$) {\pgfplotslegendfromname{legPerRepo}};
\end{tikzpicture}
\caption{Per-repository P95 merge latency and CI compute cost across 600 real OSS PRs ($5\times$ velocity, $\phi=5\%$, $80\%$ graph recall, 20 paired seeds per repo; error bars show 95\% bootstrap CIs). SSMQ achieves the lowest tail latency and runner cost on every repository.}
\label{fig:per_repo}
\end{figure*}

% AUTO-GENERATED FROM v4_exp2_main_bench — DO NOT EDIT
\begin{table}[t]
\centering
\footnotesize
\setlength{\tabcolsep}{3.0pt}
\caption{Per-Repository Breakdown at $5\times$ Velocity ($\phi=5\%$, $\rho=80\%$, 20 Seeds; Units: PR/h, min, runner-min/PR, \%).}
\label{tab:per_repo}
\begin{tabular}{llcccc}
\toprule
\textbf{Repo} & \textbf{Strategy} & \textbf{Tput} & \textbf{P95} & \textbf{Cost} & \textbf{FR\%} \\
\midrule
k8s & Static & 18.72 & 87.5 & 15.89 & 2.21\% \\
k8s & DepTree & 13.40 & 230.2 & 32.96 & 2.50\% \\
k8s & FlakeBatch & 15.41 & 162.5 & 22.48 & 0.18\% \\
k8s & \textbf{SSMQ} & \textbf{19.76} & \textbf{68.9} & \textbf{15.82} & \textbf{0.66\%} \\
\midrule
TS & Static & 13.97 & 156.1 & 23.01 & 2.10\% \\
TS & DepTree & 9.04 & 376.2 & 45.56 & 2.10\% \\
TS & FlakeBatch & 10.14 & 311.8 & 36.03 & 0.07\% \\
TS & \textbf{SSMQ} & \textbf{15.85} & \textbf{102.9} & \textbf{20.10} & \textbf{1.88\%} \\
\midrule
next.js & Static & 18.08 & 97.7 & 15.87 & 1.73\% \\
next.js & DepTree & 12.81 & 244.1 & 32.31 & 1.99\% \\
next.js & FlakeBatch & 14.78 & 173.2 & 23.52 & 0.11\% \\
next.js & \textbf{SSMQ} & \textbf{18.92} & \textbf{71.0} & \textbf{14.34} & \textbf{0.98\%} \\
\bottomrule
\end{tabular}
\end{table}


\textbf{Per-Repository Consistency and $10\times$ Velocity Scaling.} Figure~\ref{fig:per_repo} and Table~\ref{tab:per_repo} confirm that SSMQ delivers the lowest P95 latency and CI compute cost on \emph{every individual repository}, with non-overlapping $95\%$ bootstrap CIs against all baselines---including the dense \texttt{TypeScript} compiler monorepo where disjoint-4 feasibility is lowest. Under extreme $10\times$ arrival surges (Figure~\ref{fig:recall_velocity}(b), \texttt{v3\_exp3\_velocity}), SSMQ sustains \SSMQTenXThroughput{}~PR/h (\SSMQTenXPNinetyFiveLat{}~min P95) versus \StaticTenXThroughput{}~PR/h (\StaticTenXPNinetyFiveLat{}~min) for \texttt{StaticBatchBisect} and \UberTenXThroughput{}~PR/h for \texttt{UberSubmitQueue}, while \texttt{SpeculativeLinear} saturates near 5.4~PR/h due to cascading train aborts.

% AUTO-GENERATED FROM v3_exp5_graph_recall & v3_exp3_velocity — DO NOT EDIT
\begin{figure*}[t]
\centering
\begin{tikzpicture}
\begin{groupplot}[
  group style={group size=2 by 1, horizontal sep=1.5cm},
  width=0.47\textwidth, height=4.7cm,
  ymajorgrids=true, xmajorgrids=true, grid style={dashed,gray!30},
  tick label style={font=\footnotesize},
  label style={font=\footnotesize},
  title style={font=\footnotesize\bfseries},
  every axis plot/.append style={thick, mark size=2.2pt},
]
\nextgroupplot[title={(a) Non-Circularity: Off-Graph Latent Conflicts},
  xlabel={Observed Graph Recall $\rho$ (\%)}, ylabel={P95 Latency (min) $\downarrow$},
  xtick={40,60,80,100}, xmin=35, xmax=105, ymin=0, ymax=480,
  legend to name=legRecVel, legend columns=5, legend style={font=\footnotesize, draw=none, /tikz/every even column/.style={column sep=6pt}}]
\addplot[color=ssmqBlue, mark=square*] coordinates { (40,111.06) (60,93.88) (80,79.47) (100,69.29) }; \addlegendentry{SSMQ (Ours)}
\addplot[color=ssmqGray, mark=triangle*, dashed] coordinates { (40,148.65) (60,135.19) (80,121.89) (100,112.14) }; \addlegendentry{StaticBatch}
\addplot[color=ssmqOrange, mark=diamond*, dotted] coordinates { (40,442.26) (60,388.00) (80,342.68) (100,307.32) }; \addlegendentry{UberTree}
\addplot[color=ssmqTeal, mark=o, dashdotted] coordinates { (40,49.50) (100,49.50) }; \addlegendentry{SSMQ ($P_c=0$)}
\addplot[color=black!60, mark=none, densely dotted] coordinates { (40,43.73) (100,43.73) }; \addlegendentry{Static ($P_c=0$)}
\nextgroupplot[title={(b) Arrival Velocity Scaling ($1\times$--$10\times$)},
  xlabel={Offered Velocity Multiplier ($\times$)}, ylabel={Merge Throughput (PRs/hr) $\uparrow$},
  xtick={1,3,5,10}, xmin=0.5, xmax=10.5, ymin=0, ymax=28]
\addplot[color=ssmqBlue, mark=square*] coordinates { (1,4.44) (3,12.57) (5,17.78) (10,24.10) };
\addplot[color=ssmqGray, mark=triangle*, dashed] coordinates { (1,4.39) (3,11.84) (5,16.05) (10,21.34) };
\addplot[color=ssmqOrange, mark=diamond*, dotted] coordinates { (1,4.37) (3,8.55) (5,10.27) (10,10.97) };
\addplot[color=ssmqPurple, mark=pentagon*] coordinates { (1,4.52) (3,10.06) (5,12.57) (10,14.79) };
\addplot[color=ssmqRed, mark=x, dashdotted] coordinates { (1,3.50) (3,5.26) (5,5.37) (10,5.39) };
\node[font=\scriptsize, anchor=north west, fill=white, fill opacity=0.85, text opacity=1, draw=gray!40, inner sep=2pt] at (axis cs:0.8,26.8) {Purple: FlakeBatch \quad Red: SpecLinear};
\end{groupplot}
\node at ($(group c1r1.north)!0.5!(group c2r1.north)+(0,0.85cm)$) {\pgfplotslegendfromname{legRecVel}};
\end{tikzpicture}
\caption{(a) Non-circularity verification under decoupled off-graph latent conflicts ($40\%$--$100\%$ graph recall) and zero-conflict control ($P_c=0$): SSMQ retains substantial latency reductions even when $60\%$ of conflicts lie outside $\hat{G}$. (b) Throughput scaling under $1\times$--$10\times$ PR arrival bursts.}
\label{fig:recall_velocity}
\end{figure*}


\subsection{RQ3: Non-Circularity Under Off-Graph Latent Conflicts}
Table~\ref{tab:graph_recall} and Figure~\ref{fig:recall_velocity}(a) directly test whether SSMQ depends circularly on a complete dependency graph (\texttt{v3\_exp5\_graph\_recall}). When \textbf{$60\%$ of all semantic conflicts are latent off-graph edges} ($\rho=40\%$ graph recall), SSMQ still reduces P95 merge latency by over 37~minutes relative to FIFO static batching (\RecallFortySSMQPNinetyFive{} vs.\ \RecallFortyStaticPNinetyFive{}~min) and by nearly $4\times$ relative to \texttt{UberSubmitQueue}. At $\rho=60\%$ recall, SSMQ holds P95 latency to \RecallSixtySSMQPNinetyFive{} vs.\ \RecallSixtyStaticPNinetyFive{}~min.

% AUTO-GENERATED FROM v2_exp5_graph_recall — DO NOT EDIT
\begin{table}[t]
\centering
\caption{Non-Circularity Test: Off-Graph Latent Conflicts (Graph Recall $40\%$--$100\%$) and No-Conflict Control.}
\label{tab:graph_recall}
\begin{tabular}{lccc}
\toprule
\textbf{Conflict Regime} & \textbf{Static P95} & \textbf{UberTree P95} & \textbf{SSMQ P95} \\
\midrule
No-Conflict Control ($P_c=0$) & 43.7 & 153.3 & \textbf{49.5} \\
40\% Graph Recall (60\% latent) & 148.6 & 442.3 & \textbf{111.2} \\
60\% Graph Recall (40\% latent) & 135.2 & 388.0 & \textbf{93.9} \\
80\% Graph Recall (20\% latent) & 121.9 & 342.7 & \textbf{79.5} \\
100\% Graph Recall (0\% latent) & 112.1 & 307.3 & \textbf{69.4} \\
\bottomrule
\end{tabular}
\end{table}


In the \emph{Zero-Conflict Control} ($P_c = 0$), \texttt{StaticBatchBisect} achieves slightly lower P95 latency (\NoConfStaticPNinetyFive{} vs.\ \NoConfSSMQPNinetyFive{}~min) because SSMQ's protected-lane capping ($K \le 2$ on high-risk PRs) and singleton flaky guards introduce modest overhead when pairwise semantic conflicts do not exist. This clean crossover proves non-circularity and precisely delineates the operating boundary: semantic scheduling pays off whenever cross-PR module conflicts are non-trivial, with minimal tax in conflict-free workloads.

\subsection{RQ4: Component Ablation and Flaky Pareto Frontier}
Table~\ref{tab:ablation} isolates each SSMQ mechanism at the primary operating point (\texttt{v3\_exp4\_ablation}):
\begin{enumerate}
    \item \textbf{Graph-Disjoint Batching Is the Dominant Lever}: Disabling graph-disjoint risk lanes (\texttt{SSMQ-NoGraph}) degrades P95 latency from \SSMQPNinetyFiveLat{} to \AblNoGraphLat{}~min and increases CI cost from \SSMQCICost{} to \AblNoGraphCost{}~min/PR, confirming that preventing intra-batch collisions proactively drives the majority of SSMQ's speedup.
    \item \textbf{Two-Regime Flaky Gating Governs the Latency--Acceptance Frontier}: Removing flaky gating (\texttt{SSMQ-NoFlaky}) achieves lower survivor P95 latency (\AblNoFlakyLat{}~min) and compute (\AblNoFlakyCost{}~min/PR) at $\phi=5\%$ \emph{only by falsely rejecting \AblNoFlakyFRPct{}\% of innocent PRs} (vs.\ \SSMQFalseRejPct{}\% for SSMQ). Figure~\ref{fig:flaky_pareto} reveals why survivor-only evaluation is misleading: as suite flakiness rises to $\phi=10\%$ and $\phi=15\%$ (\texttt{v3\_exp6\_flaky\_sweep}), \texttt{SSMQ-NoFlaky} falsely rejects \FlakyTenNoFlakyFRPct{}\% and \FlakyFifteenNoFlakyFRPct{}\% of clean PRs (and \texttt{StaticBatchBisect} rejects \FlakyTenStaticFRPct{}\% at $\phi=10\%$). SSMQ cuts clean-PR false rejections to \FlakyTenSSMQFRPct{}\% at $\phi=10\%$ (\FlakyTenSSMQCost{}~min/PR vs.\ \FlakyTenFlakeBatchCost{}~min/PR for \texttt{FlakeAwareBatch}) and \FlakyFifteenSSMQFRPct{}\% at $\phi=15\%$ (\FlakyFifteenSSMQPNinetyFive{} vs.\ \FlakyFifteenFlakeBatchPNinetyFive{}~min P95), occupying the exact knee of the latency--acceptance Pareto frontier (Figure~\ref{fig:flaky_pareto}(b)).
    \item \textbf{1-Step Causal Blame Provides Modest Nominal Gain}: At the primary operating point, removing causal blame (\texttt{SSMQ-NoBlame}) nudges P95 latency from \SSMQPNinetyFiveLat{} to \AblNoBlameLat{}~min and cost from \SSMQCICost{} to \AblNoBlameCost{}~min/PR. We investigate this exact mechanism boundary next in RQ5.
\end{enumerate}

% AUTO-GENERATED FROM v4_exp4_ablation — DO NOT EDIT
\begin{table}[t]
\centering
\footnotesize
\setlength{\tabcolsep}{3.0pt}
\caption{Deconfounded Component Ablation ($5\times$, $\phi=5\%$, $\rho=80\%$, 20 Seeds $\times$ 3 Repos; RootSz = Mean Root Batch Size).}
\label{tab:ablation}
\begin{tabular}{lccccc}
\toprule
\textbf{Variant} & \textbf{Tput} & \textbf{P95} & \textbf{Cost} & \textbf{FR\%} & \textbf{RootSz} \\
\midrule
\textbf{SSMQ} & \textbf{18.18} & \textbf{81.0} & \textbf{16.75} & \textbf{1.17\%} & \textbf{1.44} \\
SSMQ-NoCone & 17.15 & 102.2 & 17.98 & 0.92\% & 2.09 \\
SSMQ-NoRisk & 18.09 & 82.3 & 16.80 & 1.11\% & 1.45 \\
SSMQ-NoGraph & 16.41 & 125.7 & 19.47 & 0.79\% & 2.80 \\
SSMQ-NoFlaky & 18.99 & 65.4 & 15.06 & 2.63\% & 1.43 \\
SSMQ-FixedRetry & 16.38 & 113.8 & 19.76 & 0.13\% & 1.45 \\
SSMQ-NoBlame & 18.09 & 81.8 & 16.81 & 1.21\% & 1.44 \\
StaticBatch & 16.92 & 113.8 & 18.25 & 2.01\% & 2.72 \\
\bottomrule
\end{tabular}
\end{table}

% AUTO-GENERATED FROM v3_exp6_flaky_sweep — DO NOT EDIT
\begin{figure*}[t]
\centering
\begin{tikzpicture}
\begin{groupplot}[
  group style={group size=2 by 1, horizontal sep=1.5cm},
  width=0.47\textwidth, height=4.7cm,
  ymajorgrids=true, xmajorgrids=true, grid style={dashed,gray!30},
  tick label style={font=\footnotesize},
  label style={font=\footnotesize},
  title style={font=\footnotesize\bfseries},
  every axis plot/.append style={thick, mark size=2.3pt},
]
\nextgroupplot[title={(a) Clean-PR False Rejection vs.\ Suite Flakiness},
  xlabel={Transient Flakiness Rate $\phi$ (\%)}, ylabel={Clean-PR False Rejection (\%) $\downarrow$},
  xtick={0,2,5,10,15}, xmin=-0.5, xmax=15.8, ymin=0, ymax=10.5,
  legend to name=legFlaky, legend columns=4, legend style={font=\footnotesize, draw=none, /tikz/every even column/.style={column sep=8pt}}]
\addplot[color=ssmqBlue, mark=square*] coordinates { (0,0.00) (2,0.77) (5,1.64) (10,2.23) (15,2.73) }; \addlegendentry{SSMQ (Ours)}
\addplot[color=ssmqRed, mark=triangle*, dashed] coordinates { (0,0.00) (2,0.93) (5,2.56) (10,5.34) (15,8.05) }; \addlegendentry{SSMQ-NoFlaky}
\addplot[color=ssmqGray, mark=o, dotted] coordinates { (0,0.00) (2,0.76) (5,1.96) (10,4.43) (15,6.78) }; \addlegendentry{StaticBatch}
\addplot[color=ssmqPurple, mark=diamond*, dashdotted] coordinates { (0,0.00) (2,0.00) (5,0.12) (10,0.43) (15,0.95) }; \addlegendentry{FlakeBatch}
\nextgroupplot[title={(b) Latency--Acceptance Pareto Frontier ($\phi=10\%$)},
  xlabel={P95 Merge Latency (min) $\downarrow$}, ylabel={Clean-PR False Rejection (\%) $\downarrow$},
  xmin=50, xmax=350, ymin=-0.2, ymax=6.8]
\addplot[color=ssmqBlue!45, dashed, thin, forget plot] coordinates { (70.1,5.34) (96.4,2.23) (238.8,0.43) };
\addplot[only marks, color=ssmqRed, mark=triangle*, mark size=3.2pt] coordinates { (70.1,5.34) };
\node[font=\scriptsize\bfseries, anchor=south west, inner sep=3pt, text=ssmqRed!85!black] at (axis cs:70.1,5.34) {NoFlaky};
\addplot[only marks, color=ssmqBlue, mark=square*, mark size=3.2pt] coordinates { (96.4,2.23) };
\node[font=\scriptsize\bfseries, anchor=south west, inner sep=3pt, text=ssmqBlue!85!black] at (axis cs:96.4,2.23) {SSMQ};
\addplot[only marks, color=ssmqGray, mark=o, mark size=3.2pt] coordinates { (122.3,4.43) };
\node[font=\scriptsize\bfseries, anchor=south, inner sep=3pt, text=ssmqGray!85!black] at (axis cs:122.3,4.43) {Static};
\addplot[only marks, color=ssmqPurple, mark=diamond*, mark size=3.2pt] coordinates { (238.8,0.43) };
\node[font=\scriptsize\bfseries, anchor=south, inner sep=3pt, text=ssmqPurple!85!black] at (axis cs:238.8,0.43) {FlakeBatch};
\node[font=\scriptsize, anchor=north east, fill=ssmqBlue!8, draw=ssmqBlue!50, inner sep=2.5pt] at (axis cs:342,6.4) {Better: Lower-Left $\swarrow$};
\end{groupplot}
\node at ($(group c1r1.north)!0.5!(group c2r1.north)+(0,0.85cm)$) {\pgfplotslegendfromname{legFlaky}};
\end{tikzpicture}
\caption{Flakiness sensitivity ($\phi\in[0\%,15\%]$) and latency--acceptance Pareto frontier at $\phi=10\%$. Two-regime Beta-Binomial gating positions SSMQ on the Pareto knee, slashing clean-PR false rejections relative to \texttt{SSMQ-NoFlaky} and \texttt{StaticBatchBisect} while avoiding the tail-latency explosion of unconditional full-batch retries (\texttt{FlakeAwareBatch}).}
\label{fig:flaky_pareto}
\end{figure*}


\subsection{RQ5: Failure-Density Regime and Multi-Culprit Boundary of Blame}
Why does \texttt{ASTStackCausalRanker}---despite \BenchASTTopOnePct{}\% Top-1 accuracy on real Python failures---yield only a \RegimeNomBlameGainPct{}\% P95 reduction ($p\RegimeNomBlameP{}$) over binary bisection (\texttt{SSMQ-NoBlame}) at the default operating point? Table~\ref{tab:blame_regime} and Figure~\ref{fig:blame_regime}(b) test Proposition~\ref{prop:complexity} directly across four failure-density regimes (\texttt{v3\_exp8\_blame\_regime}):

% AUTO-GENERATED FROM v5_exp8_blame_regime — DO NOT EDIT
\begin{table}[t]
\centering
\footnotesize
\setlength{\tabcolsep}{2.3pt}
\caption{Controlled $2\times 2$ Factorial Blame Sweep ($\rho=80\%$ Fixed): Survivor P95 (min), Mean Failed-Root Size (FailSz), Share of All Failed Roots With $\ge 2$ Modeled Culprits (Multi\%), and 4-Cell Holm $p$ (SSMQ vs.\ NoBlame).}
\label{tab:blame_regime}
\begin{tabular}{lcccccc}
\toprule
\textbf{Cell ($q, K$)} & \textbf{Static} & \textbf{NoBlame} & \textbf{SSMQ} & \textbf{FailSz} & \textbf{Multi\%} & \textbf{Holm $p$} \\
\midrule
8\%, $K{=}8$ & 95.0 & 66.0 & \textbf{66.0} & 1.29 & 3.6\% & 1.0000 \\
16\%, $K{=}8$ & 102.6 & 56.5 & \textbf{56.5} & 1.25 & 3.8\% & 1.0000 \\
8\%, $K{=}12$ & 105.8 & 66.0 & \textbf{66.0} & 1.29 & 3.6\% & 1.0000 \\
16\%, $K{=}12$ & 114.9 & 56.5 & \textbf{56.5} & 1.25 & 3.8\% & 1.0000 \\
\bottomrule
\end{tabular}
\end{table}

% AUTO-GENERATED FROM v3_exp0_blame_microbench & v3_exp8_blame_regime — DO NOT EDIT
\begin{figure*}[t]
\centering
\begin{tikzpicture}
\begin{groupplot}[
  group style={group size=2 by 1, horizontal sep=1.5cm},
  width=0.47\textwidth, height=4.7cm,
  ymajorgrids=true, grid style={dashed,gray!30},
  tick label style={font=\footnotesize},
  label style={font=\footnotesize},
  title style={font=\footnotesize\bfseries},
]
\nextgroupplot[title={(a) Real Python Blame Benchmark (64 Batches)},
  ybar=2pt, /pgf/bar width=10pt, ymin=0, ymax=112,
  symbolic x coords={Random,Churn,FileOverlap,ASTStack},
  xtick=data, ylabel={Accuracy (\%) $\uparrow$}, enlarge x limits=0.22,
  legend style={at={(0.03,0.97)}, anchor=north west, font=\scriptsize, draw=gray!40}]
\addplot[fill=ssmqBlue!80, draw=ssmqBlue!90!black] coordinates { (Random,26.6) (Churn,0.0) (FileOverlap,68.8) (ASTStack,93.8) }; \addlegendentry{Top-1 Acc.}
\addplot[fill=ssmqTeal!70, draw=ssmqTeal!90!black] coordinates { (Random,48.4) (Churn,9.4) (FileOverlap,100.0) (ASTStack,100.0) }; \addlegendentry{Top-2 Acc.}
\nextgroupplot[title={(b) Failure-Density Regime: Blame vs.\ Bisection},
  ybar=2pt, /pgf/bar width=8pt, ymin=0, ymax=210,
  symbolic x coords={8\%/K8,15\%/K8,22\%/K12,40\%Rec},
  xtick=data, ylabel={P95 Latency (min) $\downarrow$}, enlarge x limits=0.20,
  legend style={at={(0.03,0.97)}, anchor=north west, font=\scriptsize, draw=gray!40}]
\addplot[fill=ssmqGray!75, draw=ssmqGray!90!black] coordinates { (8\%/K8,121.9) (15\%/K8,132.3) (22\%/K12,159.7) (40\%Rec,166.6) }; \addlegendentry{StaticBatch}
\addplot[fill=ssmqOrange!80, draw=ssmqOrange!90!black] coordinates { (8\%/K8,80.6) (15\%/K8,75.7) (22\%/K12,86.0) (40\%Rec,95.2) }; \addlegendentry{SSMQ-NoBlame}
\addplot[fill=ssmqBlue!80, draw=ssmqBlue!90!black] coordinates { (8\%/K8,79.5) (15\%/K8,73.8) (22\%/K12,85.6) (40\%Rec,94.0) }; \addlegendentry{SSMQ (Full)}
\end{groupplot}
\end{tikzpicture}
\caption{(a) Attribution accuracy across 64 real executed Python failing batches, demonstrating AST symbol-level + 1-hop call-graph gains over file-level traceback matching (\texttt{FileOverlap}). (b) Failure-density regime sweep contrasting \texttt{StaticBatch}, \texttt{SSMQ-NoBlame}, and full \texttt{SSMQ}: 1-step blame yields statistically significant P95 reductions in single-culprit regimes ($8\%$--$15\%$ defect, $K=8$) while saturating when high defect density ($22\%$, $K=12$) causes multi-culprit batches that fall back to bisection.}
\label{fig:blame_regime}
\end{figure*}


\begin{itemize}
    \item \textbf{Single-Culprit Sweet Spot ($q=8\%$--$15\%, K=8$)}: Because graph-disjoint packing already eliminates most on-graph conflicts and isolates high-risk PRs into protected lanes of size $K \le 2$ (where 1-step blame and binary split are identical), batch failures primarily occur when a single latent defect strikes a fast-lane batch. As individual defect rate increases from $8\%$ to $15\%$ at $K=8$, 1-step blame's end-to-end P95 advantage over \texttt{SSMQ-NoBlame} strengthens from -\RegimeNomBlameGainPct{}\% ($p\RegimeNomBlameP{}$) to -\RegimeElevBlameGainPct{}\% (\RegimeElevSSMQPNinetyFive{} vs.\ \RegimeElevNoBlamePNinetyFive{}~min, $p\RegimeElevBlameP{}$; and nearly $2\times$ faster than \texttt{StaticBatchBisect}'s \RegimeElevStaticPNinetyFive{}~min).
    \item \textbf{Multi-Culprit Saturation ($q=22\%, K=12$)}: Exactly as predicted by Proposition~\ref{prop:complexity}, when defect density climbs to $22\%$ with $K_{\max}=12$, the expected number of simultaneous culprits inside a failing fast-lane batch exceeds two. Evicting only Top-1 suspect $s_1$ leaves a second culprit in $B \setminus \{s_1\}$, triggering binary-bisection fallback and narrowing the P95 delta to -\RegimeHighDefBlameGainPct{}\% (\RegimeHighDefSSMQPNinetyFive{} vs.\ \RegimeHighDefNoBlamePNinetyFive{}~min, $p\RegimeHighDefBlameP{}$; recovery rounds \RegimeHighDefSSMQRounds{} vs.\ \RegimeHighDefNoBlameRounds{}). Across all regimes, however, SSMQ's combined architecture outperforms FIFO \texttt{StaticBatchBisect} by over $40\%$--$46\%$, and the blame accuracy sweep (\texttt{v3\_exp7\_blame\_sweep}, $50\%$--\BenchASTTopOnePct{}\%) confirms zero safety degradation under lower ranker precision.
\end{itemize}
